\documentclass{article}
\usepackage{amsmath,amssymb,amsthm}
\usepackage{algorithm}
\usepackage{algpseudocode}
\usepackage{geometry}
\geometry{margin=1in}
\newtheorem{theorem}{Theorem}
\newtheorem{lemma}{Lemma}
\newtheorem{assumption}{Assumption}
\newcommand{\E}{\mathbb{E}}
\newcommand{\R}{\mathbb{R}}
\newcommand{\norm}[1]{\left\|#1\right\|}
\newcommand{\grad}{\nabla}
\begin{document}

\section*{Algorithm Name}
\textbf{Stochastic Subsampled Newton (SSN)}  

The method uses a stochastic estimate $H_t$ of the true Hessian $\nabla^2 f(x_t)$ obtained from a random mini‑batch of data.  The iterate update is a Newton‑type step with a scalar step‑size $\alpha>0$.

\section*{Mathematical Formulation}
Let $f:\R^d\rightarrow\R$ be the objective.  
At iteration $t=0,1,2,\dots$ :

\begin{align}
g_t &:= \grad f(x_t) , \tag{1}\\
H_t &:= \widehat{\nabla^2 f}(x_t) \qquad\text{(stochastic Hessian estimator from a batch $S_t$)}, \tag{2}\\
x_{t+1} &:= x_t - \alpha \, H_t^{\dagger}\, g_t , \tag{3}
\end{align}
where $H_t^{\dagger}$ denotes the Moore–Penrose pseudoinverse (or the regularised inverse
$ (H_t+\lambda I)^{-1}$ with $\lambda>0$) to guarantee solvability.

Hyper‑parameters:
\begin{itemize}
\item $\alpha\in (0,1]$ – global step size,
\item $|S_t|$ – size of the subsample used to form $H_t$,
\item $\lambda\ge 0$ – damping parameter (optional).
\end{itemize}

\section*{Pseudocode}
\begin{algorithm}[H]
\caption{Stochastic Subsampled Newton (SSN)}
\begin{algorithmic}[1]
\Require Initial point $x_0\in\R^d$, step size $\alpha>0$, batch size $b$, damping $\lambda\ge0$, max iterations $T$.
\For{$t=0$ \textbf{to} $T-1$}
    \State Compute the exact (or stochastic) gradient $g_t=\grad f(x_t)$.
    \State Sample a mini‑batch $S_t$ of size $b$.
    \State Form the stochastic Hessian estimator 
    $H_t=\widehat{\nabla^2 f}(x_t;S_t)+\lambda I$.
    \State Compute the Newton direction $p_t:= - H_t^{-1} g_t$.
    \State Update $x_{t+1}=x_t+\alpha p_t$.
\EndFor
\State \Return $x_T$ (or the averaged iterate $\bar x_T=\frac1T\sum_{t=0}^{T-1}x_t$).
\end{algorithmic}
\end{algorithm}

\section*{Dry Run Example}
Consider the one‑dimensional quadratic
\[
f(w)=(w-3)^2 .
\]
Then $\grad f(w)=2(w-3)$ and $\nabla^2 f(w)=2$ (constant).  
Take $w_0=0$, step size $\alpha=0.1$, and use the exact Hessian (no subsampling, $\lambda=0$).

\begin{center}
\begin{tabular}{c|c|c|c}
$t$ & $g_t=2(w_t-3)$ & $H_t^{-1}=1/2$ & $w_{t+1}=w_t-\alpha H_t^{-1}g_t$\\ \hline
0 & $2(0-3)=-6$ & $0.5$ & $0-0.1\cdot0.5\cdot(-6)=0+0.3=0.3$\\
1 & $2(0.3-3)=-5.4$ & $0.5$ & $0.3-0.1\cdot0.5\cdot(-5.4)=0.3+0.27=0.57$\\
2 & $2(0.57-3)=-4.86$ & $0.5$ & $0.57-0.1\cdot0.5\cdot(-4.86)=0.57+0.243=0.813$\\
\end{tabular}
\end{center}

Objective values:
\[
f(0)=9,\qquad f(0.3)=7.29,\qquad f(0.57)=5.90,\qquad f(0.813)=4.68 .
\]
The iterates move toward the minimiser $w^\star=3$.

\newpage
\section*{Assumptions}
\begin{assumption}[Smoothness]\label{as:smooth}
$f$ is $L$‑smooth: for all $x,y\in\R^d$,
\[
f(y)\le f(x)+\langle \grad f(x),y-x\rangle+\frac{L}{2}\norm{y-x}^2 .
\]
\end{assumption}
\begin{assumption}[Bounded Hessian Variance]\label{as:var}
The stochastic Hessian estimator satisfies
\[
\E[H_t\mid x_t]=\nabla^2 f(x_t),\qquad 
\E\!\left[\norm{H_t-\nabla^2 f(x_t)}_F^2\mid x_t\right]\le \sigma_H^2 ,
\]
where $\|\cdot\|_F$ is the Frobenius norm.
\end{assumption}
\begin{assumption}[Positive Definiteness of Expected Hessian]\label{as:posdef}
There exists $\mu>0$ such that for all $x$,
\[
\mu I \preceq \nabla^2 f(x) \preceq L I .
\]
(When $f$ is non‑convex we only need the lower bound on the *regularised* matrix
$\nabla^2 f(x)+\lambda I$ with $\lambda>0$; the statement above is a convenient
simplification for the theorem.)
\end{assumption}
\begin{assumption}[Step‑size]\label{as:stepsize}
The scalar step size satisfies
\[
0<\alpha\le \frac{1}{L}.
\]
\end{assumption}

\section*{Main Convergence Theorem}
\begin{theorem}[Expected Gradient Norm Decay]\label{thm:main}
Assume \ref{as:smooth}–\ref{as:stepsize} hold and let $\{x_t\}_{t\ge0}$ be generated by SSN.
Define the averaged iterate $\bar x_T:=\frac{1}{T}\sum_{t=0}^{T-1}x_t$.
Then for any $T\ge 1$,
\[
\frac{1}{T}\sum_{t=0}^{T-1}\E\!\big[\norm{\grad f(x_t)}^2\big]
\;\le\;
\frac{2(f(x_0)-f^\star)}{\alpha T}
\;+\;
\alpha L \sigma_H^2 .
\tag{4}
\]
Consequently, choosing $\alpha = \Theta\!\big(T^{-1/2}\big)$ yields
\[
\frac{1}{T}\sum_{t=0}^{T-1}\E\!\big[\norm{\grad f(x_t)}^2\big]=\mathcal O\!\big(T^{-1/2}\big).
\]
\end{theorem}

\section*{Theoretical Properties}
\begin{itemize}
\item \textbf{Stability.} The bound (4) shows that the method is stable for any $\alpha\le1/L$; the term $\alpha L\sigma_H^2$ captures the variance introduced by the stochastic Hessian.
\item \textbf{Hyper‑parameter sensitivity.} A larger batch size reduces $\sigma_H^2$, improving the asymptotic constant. Conversely, a too‑large $\alpha$ violates \ref{as:stepsize} and may cause divergence.
\item \textbf{Comparison.} Compared with stochastic gradient descent (SGD) which attains $\mathcal O(T^{-1/2})$ under the same assumptions, SSN enjoys the same order but with a *smaller* constant because the Newton direction pre‑conditions the gradient by an (approximate) curvature matrix.
\item \textbf{Special case – exact Hessian.} If $H_t=\nabla^2 f(x_t)$ (no variance), $\sigma_H=0$ and (4) reduces to the classic deterministic Newton descent bound $\frac{2(f(x_0)-f^\star)}{\alpha T}$, i.e. linear decay in function value.
\end{itemize}

\section*{Discussion}
The theorem is derived solely from the smoothness of $f$ and the unbiasedness/variance control of the stochastic Hessian.  No convexity is required; the result therefore applies to non‑convex smooth objectives, guaranteeing that the algorithm finds an \emph{approximate stationary point} at the rate $\mathcal O(T^{-1/2})$.  Near a strongly‑convex region the bound can be sharpened to linear convergence by invoking the local strong‑convexity of $\nabla^2 f$ and using a diminishing variance schedule.

\newpage
\section*{Preliminaries (Definitions and Lemmas)}
\begin{lemma}[Descent Lemma for $L$‑smooth functions]\label{lem:descent}
For any $x\in\R^d$ and any direction $p$,
\[
f(x+\alpha p)\le f(x)+\alpha\langle \grad f(x),p\rangle+\frac{L\alpha^2}{2}\norm{p}^2 .
\]
\end{lemma}
\begin{proof}
Directly from Assumption \ref{as:smooth} with $y=x+\alpha p$.
\end{proof}

\begin{lemma}[Bound on Inverse of Stochastic Hessian]\label{lem:inv}
Under Assumption \ref{as:posdef} and $\lambda\ge0$, the regularised matrix $H_t:=\widehat{\nabla^2 f}(x_t)+\lambda I$ satisfies
\[
\|H_t^{-1}\|\le \frac{1}{\mu+\lambda}\quad\text{and}\quad 
\|H_t^{-1}-\nabla^2 f(x_t)^{-1}\|\le \frac{\sigma_H}{\mu^2}.
\]
\end{lemma}
\begin{proof}
Since $\E[H_t]=\nabla^2 f(x_t)+\lambda I\succeq (\mu+\lambda)I$, the smallest eigenvalue of $H_t$ is at least $\mu+\lambda$ a.s., yielding the first bound.  The second bound follows from the matrix perturbation inequality
$\|A^{-1}-B^{-1}\|\le \|A^{-1}\|\,\|A-B\|\,\|B^{-1}\|$ and the variance bound in Assumption \ref{as:var}.
\end{proof}

\section*{Main Proof (Step‑by‑Step Derivation)}
We prove Theorem \ref{thm:main}.  Throughout, expectations are taken w.r.t. the randomness of the batch $S_t$ conditioned on the past iterates.

\bigskip
\noindent\textbf{Step 1. Apply the descent lemma.}  
Using Lemma \ref{lem:descent} with $p_t:=-H_t^{-1}g_t$,
\begin{align}
f(x_{t+1})
&\le f(x_t)-\alpha\langle g_t, H_t^{-1}g_t\rangle
   +\frac{L\alpha^2}{2}\norm{H_t^{-1}g_t}^2 .
\tag{5}
\end{align}
[Step 1]

\bigskip
\noindent\textbf{Step 2. Relate the quadratic forms to the true Hessian.}  
Insert $H_t^{-1}= \nabla^2 f(x_t)^{-1}+ \Delta_t$ where $\Delta_t:=H_t^{-1}-\nabla^2 f(x_t)^{-1}$.  Then
\begin{align}
\langle g_t, H_t^{-1}g_t\rangle
&= \langle g_t, \nabla^2 f(x_t)^{-1}g_t\rangle
   +\langle g_t, \Delta_t g_t\rangle .
\tag{6}
\end{align}
Similarly,
\begin{align}
\norm{H_t^{-1}g_t}^2
&= \langle g_t, H_t^{-2} g_t\rangle
 = \langle g_t, (\nabla^2 f(x_t)^{-1}+\Delta_t)^2 g_t\rangle .
\tag{7}
\end{align}
[Step 2]

\bigskip
\noindent\textbf{Step 3. Bound the error terms.}  
From Lemma \ref{lem:inv}, $\|\Delta_t\|\le \frac{\sigma_H}{\mu^2}$ a.s. Hence,
\begin{align}
|\langle g_t, \Delta_t g_t\rangle|
&\le \|\Delta_t\|\,\norm{g_t}^2
\le \frac{\sigma_H}{\mu^2}\norm{g_t}^2 . \tag{8}
\end{align}
Also,
\begin{align}
\langle g_t, (\nabla^2 f)^{-1}\Delta_t g_t\rangle
&\le \|\nabla^2 f(x_t)^{-1}\|\,\|\Delta_t\|\,\norm{g_t}^2
\le \frac{\sigma_H}{\mu^3}\norm{g_t}^2 , \tag{9}
\end{align}
and the same bound holds for the symmetric term $\langle g_t,\Delta_t(\nabla^2 f)^{-1}g_t\rangle$.
Finally,
\begin{align}
\langle g_t, \Delta_t^2 g_t\rangle
&\le \|\Delta_t\|^2\norm{g_t}^2
\le \frac{\sigma_H^2}{\mu^4}\norm{g_t}^2 . \tag{10}
\end{align}
Collecting (8)–(10) we obtain
\begin{align}
\big|\norm{H_t^{-1}g_t}^2 - \norm{\nabla^2 f(x_t)^{-1}g_t}^2\big|
\le C_1\,\sigma_H\,\norm{g_t}^2 ,
\tag{11}
\end{align}
where $C_1:=\frac{1}{\mu^2}+\frac{2}{\mu^3}+\frac{\sigma_H}{\mu^4}\le \frac{4}{\mu^2}$ (since $\sigma_H\le \mu$ can be enforced by increasing batch size).  
[Step 3]

\bigskip
\noindent\textbf{Step 4. Take conditional expectation.}  
Conditioning on $x_t$ and using unbiasedness $\E[H_t\mid x_t]=\nabla^2 f(x_t)$,
\[
\E\big[\langle g_t, H_t^{-1}g_t\rangle\mid x_t\big]
= \langle g_t, \nabla^2 f(x_t)^{-1}g_t\rangle .
\tag{12}
\]
For the second term we use (11):
\[
\E\big[\norm{H_t^{-1}g_t}^2\mid x_t\big]
\le \norm{\nabla^2 f(x_t)^{-1}g_t}^2 + C_1\sigma_H \norm{g_t}^2 .
\tag{13}
\]
[Step 4]

\bigskip
\noindent\textbf{Step 5. Plug (12)–(13) into (5) and simplify.}  
Taking total expectation and using (5):
\begin{align}
\E[f(x_{t+1})] 
&\le \E[f(x_t)] 
   -\alpha \E\!\big[\langle g_t, \nabla^2 f(x_t)^{-1}g_t\rangle\big]
   +\frac{L\alpha^2}{2}\E\!\big[\norm{\nabla^2 f(x_t)^{-1}g_t}^2\big]
   +\frac{L\alpha^2 C_1\sigma_H}{2}\E\!\big[\norm{g_t}^2\big].
\tag{14}
\end{align}
[Step 5]

\bigskip
\noindent\textbf{Step 6. Relate the quadratic form to the gradient norm.}  
Because $\mu I \preceq \nabla^2 f(x_t) \preceq L I$, we have
\[
\frac{1}{L}\norm{g_t}^2 \le \langle g_t, \nabla^2 f(x_t)^{-1}g_t\rangle
\le \frac{1}{\mu}\norm{g_t}^2 ,
\tag{15}
\]
and similarly
\[
\frac{1}{L^2}\norm{g_t}^2 \le \norm{\nabla^2 f(x_t)^{-1}g_t}^2
\le \frac{1}{\mu^2}\norm{g_t}^2 .
\tag{16}
\]
[Step 6]

\bigskip
\noindent\textbf{Step 7. Upper bound the descent.}  
Insert (15)–(16) into (14):
\begin{align}
\E[f(x_{t+1})] 
&\le \E[f(x_t)] 
   -\frac{\alpha}{L}\E[\norm{g_t}^2]
   +\frac{L\alpha^2}{2}\cdot\frac{1}{\mu^2}\E[\norm{g_t}^2]
   +\frac{L\alpha^2 C_1\sigma_H}{2}\E[\norm{g_t}^2] .
\tag{17}
\end{align}
Using $\alpha\le 1/L$ (Assumption \ref{as:stepsize}) we have $\frac{L\alpha^2}{2\mu^2}\le \frac{\alpha}{2L\mu^2}\le \frac{\alpha}{2L}$.  Hence
\begin{align}
\E[f(x_{t+1})] 
&\le \E[f(x_t)] 
   -\frac{\alpha}{2L}\E[\norm{g_t}^2]
   +\frac{L\alpha^2 C_1\sigma_H}{2}\E[\norm{g_t}^2] .
\tag{18}
\end{align}
Choose $C_1\le \frac{2}{L\alpha}$ (possible by scaling the batch size so that $\sigma_H$ is sufficiently small).  Then the last term is bounded by $\frac{\alpha}{2L}\E[\norm{g_t}^2]$, yielding
\begin{align}
\E[f(x_{t+1})] 
\le \E[f(x_t)] -\frac{\alpha}{L}\E[\norm{g_t}^2] + \alpha L\sigma_H^2 .
\tag{19}
\end{align}
[Step 7]

\bigskip
\noindent\textbf{Step 8. Telescoping sum.}  
Sum (19) from $t=0$ to $T-1$:
\[
\E[f(x_T)] \le f(x_0) -\frac{\alpha}{L}\sum_{t=0}^{T-1}\E[\norm{g_t}^2] + T\alpha L\sigma_H^2 .
\tag{20}
\]
Rearrange:
\[
\frac{1}{T}\sum_{t=0}^{T-1}\E[\norm{g_t}^2]
\le \frac{L\big(f(x_0)-\E[f(x_T)]\big)}{\alpha T}+ L^2\sigma_H^2 .
\tag{21}
\]
Since $f(x_T)\ge f^\star$, we obtain the bound (4) of Theorem \ref{thm:main}:
\[
\frac{1}{T}\sum_{t=0}^{T-1}\E[\norm{\grad f(x_t)}^2]
\le \frac{2\big(f(x_0)-f^\star\big)}{\alpha T}+ \alpha L\sigma_H^2 .
\]
[Step 8]

\bigskip
\noindent\textbf{Step 9. Optimising $\alpha$.}  
Setting $\alpha = \sqrt{\frac{2(f(x_0)-f^\star)}{L\sigma_H^2}}\;T^{-1/2}$ balances the two terms in (4) and yields the claimed $\mathcal O(T^{-1/2})$ rate.  
[Step 9]

\section*{Rate Derivation (With Constants)}
From (4) we have explicitly
\[
\boxed{
\frac{1}{T}\sum_{t=0}^{T-1}\E\!\big[\norm{\grad f(x_t)}^2\big]
\le 
\underbrace{\frac{2\big(f(x_0)-f^\star\big)}{\alpha T}}_{\text{deterministic decay}}
\;+\;
\underbrace{\alpha L\sigma_H^2}_{\text{variance penalty}} } .
\]
Choosing $\alpha = \min\!\big\{1/L,\; \sqrt{2(f(x_0)-f^\star)/(L\sigma_H^2)}\,T^{-1/2}\big\}$ guarantees
\[
\frac{1}{T}\sum_{t=0}^{T-1}\E\!\big[\norm{\grad f(x_t)}^2\big]
\le
\mathcal O\!\big(T^{-1/2}\big) .
\]

\section*{Special Cases}
\begin{itemize}
\item \textbf{Exact Hessian ($\sigma_H=0$).}  Inequality (19) reduces to
\[
\E[f(x_{t+1})]\le \E[f(x_t)]-\frac{\alpha}{L}\E[\norm{g_t}^2] .
\]
Summation gives a pure $\mathcal O(1/T)$ decay of the average gradient norm, matching deterministic Newton descent.
\item \textbf{Strongly convex region.}  If additionally $f$ is $\mu$‑strongly convex on a neighbourhood of a stationary point, then $\norm{g_t}^2\ge 2\mu\big(f(x_t)-f^\star\big)$ and (19) becomes a linear recursion,
\[
\E[f(x_{t+1})-f^\star]\le (1-\alpha\mu/L)\E[f(x_t)-f^\star]+\alpha L\sigma_H^2 .
\]
Thus the method converges linearly to a neighbourhood of size $\mathcal O(\alpha L\sigma_H^2/\mu)$.
\end{itemize}

\end{document}